\documentclass[12pt]{article}
\usepackage[spanish]{babel}
\usepackage[utf8]{inputenc}
\usepackage{amsmath, amssymb}
\usepackage{geometry}
\usepackage{graphicx}
\usepackage{fancyhdr}

\geometry{letterpaper, margin=2.5cm}

\pagestyle{fancy}
\fancyhf{}
\rhead{Matemáticas para Ciencia de Datos}
\lhead{Actividad 3.1}
\rfoot{\thepage}

\title{\textbf{Actividad de Aprendizaje 3.1 \\ Conceptos de Probabilidad}}
\author{Emilio Izquierdo Montero \\ Grupo 8SA}
\date{Marzo 2026}

\begin{document}

\maketitle

\section*{Introducción}
La probabilidad es una herramienta fundamental en la ciencia de datos, ya que permite modelar la incertidumbre y tomar decisiones informadas en base a sucesos anteriores, en este documento se presentan ejercicios resueltos sobre distintos conceptos clave.

\newpage

%------------------------------------------------
\section{Permutaciones y Combinaciones}

\subsection*{Concepto}
Las permutaciones consideran el orden de los elementos, mientras que las combinaciones no.

\begin{equation}
P(n,r)=\frac{n!}{(n-r)!}, \quad C(n,r)=\frac{n!}{r!(n-r)!}
\end{equation}

\subsection*{Ejercicios}

\textbf{1.} ¿De cuántas formas se pueden ordenar 6 libros distintos en un estante?

\[
P(6,6)=6!=720
\]

\textbf{2.} ¿Cuántas palabras distintas se pueden formar con las letras de la palabra ``DATOS''?

\[
5!=120
\]

\textbf{3.} Seleccionar 4 estudiantes de un grupo de 12.

\[
C(12,4)=495
\]

\textbf{4.} ¿Cuántos comités de 3 personas pueden formarse con 8 candidatos?

\[
C(8,3)=56
\]

\textbf{5.} Ordenar 3 posiciones de entre 10 personas.

\[
P(10,3)=720
\]

\textbf{Árbol de decisiones o diagrama factorial}
\begin{figure}[h]
\centering
\includegraphics[width=0.5\textwidth]{/home/emilio/Descargas/arboldedecision.jpg}
\caption{Ejemplo de representación de permutaciones}
\end{figure}

%------------------------------------------------
\section{Probabilidad Condicional}

\subsection*{Concepto}

\begin{equation}
P(A|B)=\frac{P(A \cap B)}{P(B)}
\end{equation}

\subsection*{Ejercicios}

\textbf{1.} En una baraja, ¿cuál es la probabilidad de sacar una reina dado que la carta es figura?

\[
P = \frac{4}{12} = \frac{1}{3}
\]

\textbf{2.} Se sabe que el 70\% de los estudiantes estudian y de ellos el 80\% aprueba.

\[
P(A|B)=0.8
\]

\textbf{3.} En un dado, probabilidad de obtener un número par dado que es mayor que 3.

\[
\frac{2}{3}
\]

\textbf{4.} Probabilidad de lluvia dado que está nublado (supongamos datos).

\[
P = 0.6
\]

\textbf{5.} Probabilidad de enfermedad dado que el paciente tiene síntomas.

\[
P = \frac{P(\text{síntomas} \cap \text{enfermedad})}{P(\text{síntomas})}
\]

\textbf{Diagrama de Venn}
\begin{figure}[h]
\centering
\includegraphics[width=0.5\textwidth]{/home/emilio/Descargas/venn.png}
\caption{Eventos condicionados}
\end{figure}

%------------------------------------------------
\section{Eventos Dependientes}

\subsection*{Concepto}

\[
P(A \cap B)=P(A)P(B|A)
\]

\subsection*{Ejercicios}

\textbf{1.} Sacar dos cartas rojas sin reemplazo:

\[
\frac{26}{52} \cdot \frac{25}{51}
\]

\textbf{2.} Extraer dos bolas rojas de una urna sin reemplazo.

\textbf{3.} Seleccionar dos estudiantes consecutivos.

\textbf{4.} Lanzar dos monedas dependientes.

\textbf{5.} Probabilidad de dos eventos consecutivos.

\textbf{Árbol de probabilidad}
\begin{figure}[h]
\centering
\includegraphics[width=0.5\textwidth]{/home/emilio/Descargas/dependent.png}
\caption{Eventos dependientes}
\end{figure}

%------------------------------------------------
\section{Eventos No Excluyentes}

\subsection*{Concepto}

\begin{equation}
P(A \cup B)=P(A)+P(B)-P(A \cap B)
\end{equation}

\subsection*{Ejercicios}

\textbf{1.} Probabilidad de trabajar y estudiar:

\[
0.5 + 0.4 - 0.2 = 0.7
\]

\textbf{2.} Probabilidad de números pares o múltiplos de 3.

\textbf{3.} Probabilidad de tomar café o té.

\textbf{4.} Probabilidad de aprobar dos materias.

\textbf{5.} Probabilidad de practicar dos deportes.

\textbf{Diagrama de Venn}
\begin{figure}[h]
\centering
\includegraphics[width=0.5\textwidth]{/home/emilio/Descargas/bdcd.jpg}
\caption{Eventos no excluyentes}
\end{figure}

%------------------------------------------------
\section{Teorema de Bayes}

\subsection*{Concepto}

\begin{equation}
P(A|B)=\frac{P(B|A)P(A)}{P(B)}
\end{equation}

\subsection*{Ejercicios}

\textbf{1.} Una prueba médica tiene 95\% de sensibilidad y 90\% de especificidad.

\textbf{2.} Probabilidad de que un producto defectuoso provenga de una máquina.

\textbf{3.} Clasificación de spam.

\textbf{4.} Detección de fraude.

\textbf{5.} Diagnóstico médico.

\textbf{Espacio para figura: Diagrama de árbol Bayesiano}
\begin{figure}[h]
\centering
\includegraphics[width=0.5\textwidth]{/home/emilio/Descargas/diagramavayes.jpg}
\caption{Aplicación del Teorema de Bayes}
\end{figure}

%------------------------------------------------
\section*{Conclusión}

Los conceptos de probabilidad son esenciales para el análisis de datos, ya que permiten modelar incertidumbre y realizar inferencias basadas en información parcial.

\end{document}
